\section{Method}
\label{sec:method}

\paragraph{Paired updates and probes.}
Let $h^{(l)}(C)\in\mathbb R^{d}$ be the residual stream after block $l$ at the anchor for context $C$. For increment $C_{k-1}\to C_k$ we form $\Delta h^{(l)}=h^{(l)}(C_k)-h^{(l)}(C_{k-1})$ and, for each task and layer, fit an $\ell_2$-regularised logistic regression on standardised features with class balancing, selecting the layer by grouped three-fold cross-validation on training questions only. We also report the \emph{difference-of-means} direction $\vdm\propto\mathbb E[\Delta h\mid y{=}1]-\mathbb E[\Delta h\mid y{=}0]$ and, for the raw state $h(C)$, \emph{state probes} for sufficiency and correctness. Directions are unit vectors in raw activation space.

\paragraph{Nuisance analysis.}
A paired difference also encodes added tokens, context length, answer-string presence, passage relevance and the change in the output distribution. We build a nuisance vector $z$ per increment (token counts, gold-answer log-probability change, next-token entropy and top-probability change, answer-string presence, TF-IDF relevance, insertion position, and a 16-dimensional TF-IDF-SVD embedding of the added text; \cref{app:probing}) and report (i) a nuisance-only classifier, (ii) the probe on residuals of a ridge regression of $\Delta h$ on $z$, and (iii) the AUROC gain of adding out-of-fold probe scores to the nuisance model.

\paragraph{Text baselines.}
Text-only classifiers on the same increments and splits: TF-IDF logistic regression and a fine-tuned DeBERTa-v3-base cross-encoder, on ``question [SEP] added paragraph'' alone and with the previous context appended.

\paragraph{Transfer.}
Learned directions are applied without retraining to (i) the \emph{conversational revision} protocol, in which the model sees its own previous answer $A_k$ and only the new evidence and is told to change its answer only if the total evidence warrants it, pairing $h_{\text{conv}}(C_k\to C_{k+1})-h(C_k)$; (ii) \twowiki{}; (iii) 3--4-hop \musique{}; and (iv) the datasets of \cref{sec:external}. Model-family transfer is replication of the procedure, since hidden sizes differ.

\paragraph{Causal tests.}
Interventions act on held-out questions of the document split. \emph{Steering} adds $\alpha\,u\,v$ from the anchor onward at layer $l$ (or all positions), where $u$ is the separation between the sufficient and insufficient state means along $v$ on training data, so $\alpha=1$ moves an activation by one sufficient-minus-insufficient gap; a random unit direction with the same $\ell_2$ magnitude is the control. \emph{Projection-out} sets the component along $v$ to its insufficient-state mean at all layers from $l$ onward. \emph{Patching} runs the insufficient prompt $C_{n-1}$ with its anchor residual replaced, at every layer from $l$ onward, by $h_{\text{suff}}$ (full), $h_{\text{ins}}+\mathrm{proj}_v(\Delta)$ (parallel), $h_{\text{ins}}+\Delta-\mathrm{proj}_v(\Delta)$ (orthogonal), the projection of $\Delta$ onto a random direction, onto the top-$k$ principal components of training deltas, or onto the gold answer's first-token unembedding direction. Outcomes are the rate of non-abstention (``answered'') and of correct answers on 150 held-out questions per set; for false-evidence contexts we also report the rate of producing the injected false answer.
